\documentclass[11pt]{article}
\usepackage[T1]{fontenc}
\usepackage{lmodern,microtype,amsmath,amssymb,amsthm,booktabs,needspace}
\usepackage[margin=0.92in]{geometry}
\usepackage[hidelinks]{hyperref}
\hypersetup{pdftitle={One-runner extension and phase stability of a 57-runner shifted configuration},
           pdfauthor={Chengyi Fang}}
\setlength{\emergencystretch}{1em}
\allowdisplaybreaks
\newtheorem{theorem}{Theorem}
\newtheorem{lemma}[theorem]{Lemma}
\newtheorem{corollary}[theorem]{Corollary}
\newtheorem{proposition}[theorem]{Proposition}
\newcommand{\T}{\mathbb{R}/\mathbb{Z}}
\newcommand{\Z}{\mathbb{Z}}
\newcommand{\norm}[1]{\left\lVert #1\right\rVert}
\DeclareMathOperator{\lcm}{lcm}
\title{One-runner extension and phase stability of a\\
       57-runner shifted configuration}
\author{Chengyi Fang\\[-2pt]\small
\href{https://orcid.org/0009-0006-1237-7763}{ORCID 0009-0006-1237-7763}}
\date{Revised 27 November 2026}

\begin{document}
\maketitle
\begin{abstract}
Fix the 57-runner configuration credited to \texttt{tomdif} in the
public shifted lonely runner certificate collection. Over every
additional distinct positive integer speed and arbitrary real phase,
the optimal covering radius is \(1359451/79000000\). A single speed,
\(528344\), attains it, with a closed interval of optimal phases.
The exact radius and a perturbation barrier are Lean-verified.
An independent weighted residual certificate shows that seven extra
runners, even with repeated speeds, cannot cover the time circle at
radius \(1/59\). Any 58-runner cover at that radius shares at most
50 speed--phase pairs with the fixed configuration. We also give
restricted exclusion computations for coordinated phase changes and
for several deletion-and-insertion problems.
\end{abstract}

\section{Setup}

Let \(\norm{x}=\min_{k\in\Z}|x-k|\). For a family \(C\) of
speed--phase pairs \((v,s)\) with \(v\in\Z_{>0}\) and \(s\in\T\), define
\[
 d_C(t)=\min_{(v,s)\in C}\norm{vt+s},\qquad
 \gamma(C)=\max_{t\in\T}d_C(t).
\]
The observer at the origin contributes no moving runner. A shifted
lonely runner counterexample with \(n\) distinct speeds satisfies
\(\gamma(C)<1/(n+1)\). Blanco, Criado and Santos give counterexamples
for \(5\le n\le17\) \cite[Theorem 1.8]{BCS}, and Poliakova
gives them for \(n\ge95\) \cite[Theorem 2]{Poliakova}. The public
collection \cite{upstream} gives exact rational certificates at each
\(n=18,\ldots,57\), since verified in Lean \cite{upstreamlean}.
Zhang's Question 8.5 asks both about all sizes \(n\ge5\) and about
global minimizers \cite{Zhang}. Here the primary results hold for one
specified seed, and the later tests use specified speed and phase
restrictions.

Let \(B\) be the 57-runner configuration from
\texttt{certificates.json} in the public collection at commit
\href{https://github.com/1337isnot1337/solved-conjectures/blob/75a0b27a1b3e76715a45704865454b14f8f1c63d/results/shifted-lonely-runner-counterexamples/certificates.json}{\texttt{75a0b27a}}.
Its coordinates are reproduced in the companion
\texttt{known\_slr\_57.json}. The exact value
\(\gamma(B)=43021/2500000\) comes from that certificate.
The public verifier checks all 40 covers in rational arithmetic and
recomputes the maximum at \(n=57\) by its full event enumeration.
We use \(d_0=1/59\) for the 58-runner target.

For time intervals at radius \(d\in(0,1/2)\), the bands of \((v,s)\)
are \([(k-s-d)/v,(k-s+d)/v]\) with \(k\in\Z\). For strict coverage
take the bands open. We work on \([0,1]\) with \(0\) and \(1\)
identified. The \(v\) bands modulo one for a single runner are
disjoint, apart from the representation at this seam.

\Needspace{9\baselineskip}
\section{Exact optimum with one extra runner}

\begin{theorem}\label{thm:frozen}
Among configurations obtained from \(B\) by adding one runner with
a distinct positive integer speed and arbitrary real phase, the
least covering radius is
\[
 d_*=\frac{1359451}{79000000}.
\]
The speed is uniquely \(528344\). At that speed the full set of
optimal phases modulo one is the closed interval with left endpoint
\(4737613/31600000\) and length \(178783/158000000\).
\end{theorem}

\begin{proof}
For weak coverage at a specified \(d\), sweep and merge the closed
bands of \(B\). The complement comprises the open gaps where
\(d_B(t)>d\). A new runner covers an open gap if and only if it
covers its closure, by continuity. For the seed here \(0\) is covered.
At \(d=d_*\) there are eleven positive-length gaps. The longest has
length \(53/1264000000\). Since the phase image of a time interval
\([a,b]\) is connected, covering it requires \(v(b-a)\le2d\).
In particular any possible speed at \(d_*\) is at most \(820800\).

For each candidate speed \(v\), order the closed gaps by decreasing
length and write the first as \([a_0,b_0]\).
The gap \([a,b]\) restricts the phase to the circular window
\[
 [ -va-d,\,-vb+d]\pmod1,
 \tag{1}\label{eq:window}
\]
empty when the endpoints are reversed. With
\(u=s+va_0+d\bmod1\), the first window sets
\(0\le u\le 2d-v(b_0-a_0)\). Every other window starts at
\(v(a_0-a)\bmod1\) and has length \(2d-v(b-a)\).
Choose each start in \([-1/2,1/2]\), and discard negative lengths.
Here \(2d<1/4\), so the range of \(u\) can intersect only this
representative, without a translate by a nonzero integer.
Intersect the closed intervals in this coordinate.

For \(d=d_*\) the finite enumeration leaves exactly one new speed
and the stated interval of phases. The arithmetic can be reproduced
without floating-point operations. Multiply \(d\) and the gap
endpoints by their least common denominator \(D\). For each
\(1\le v\le820800\) absent from the seed, initialize
\[
 (\ell,h)=(0,\,2dD-vD(b_0-a_0)).
\]
For every other gap set \(z\) to \(vD(a_0-a)\) reduced modulo \(D\)
into \([-D/2,D/2]\). Update
\[
  \ell\gets\max(\ell,z),\qquad
  h\gets\min(h,z+2dD-vD(b-a)).
\]
Drop the speed as soon as \(\ell>h\). The single surviving interval,
after dividing by \(D\) and shifting by \(-va_0-d\), is as stated.
At its left endpoint the eleven closed gaps also pass a direct
test of (1).

For the lower bound use
\[
 t_*=\frac{38586773}{79000000}.
\]
At this point all seed distances are at least \(d_*\), with equality
for speeds \(44\) and \(114\). These inequalities are checked by
substituting the rational seed. The distance of the only feasible
new speed at its leftmost phase is \(66324111/158000000\) there.
Changing the phase anywhere inside the feasible interval can
decrease the distance by at most \(178783/158000000\), so the new
distance is still strictly greater than \(d_*\). Thus no weak cover
at \(d_*\) can have smaller radius. The isolated point \(t_*\)
belongs to the complement of the \emph{open} seed bands at \(d_*\);
the lower-bound argument uses it even though the complement of the
closed bands has only positive-length gaps. Source and certificate
are in \texttt{optimal\_extension.py} and
\texttt{optimal\_extension.json}.
\end{proof}

An independent Lean proof of the radius uses six pairs of rational
times where all seed distances are at least \(d_*\).
For every \(1\le v\le820800\) at least one pair \((a,b)\) satisfies
\(\norm{v(a-b)}\ge2d_*\). The Lean kernel checks the finite modular
inequality. If the new runner had distance strictly less than \(d_*\)
at both times, the triangle inequality on \(\T\) would contradict
it. For larger \(v\) the proof uses two times separated by
\(2d_*/v\) inside the longest closed gap of the seed coverage.
The upper bound in Lean uses a chain of rational closed bands for
the attaining configuration. This proves the radius with no bound
on speeds and no assumption of rational phases. The uniqueness
and phase interval in Theorem~\ref{thm:frozen} use the Python
enumeration.

\begin{corollary}\label{cor:stability}
Keep the 57 speeds of \(B\) and move each of its phases by at most
\(\epsilon\) in circular distance. With any additional runner the
covering radius is at least \(d_*-\epsilon\). At
\(\epsilon=1207609/4661000000\) this is \(1/59\).
\end{corollary}
\begin{proof}
For fixed time and speed the distance to the nearest integer is
1-Lipschitz in the phase on \(\T\). The same Lipschitz bound, in the
maximum displacement over runners, holds after taking a minimum
over runners and a maximum over time. Compare to the old phases
with the extra runner unchanged. The value
\(d_*-1/59=1207609/4661000000\) is exact.
A common change of time origin replaces all phases by \(s_i+v_i t_0\)
without changing the maximum. The barrier holds in any box about
such a time translate of \(B\).
\end{proof}

\section{A weighted obstruction to seven additions}

\begin{theorem}\label{thm:weights}
If all runners of \(B\) are retained, appending any seven runners
of positive integer speed and arbitrary real phase cannot achieve
radius at most \(1/59\). Repeated speeds among the extra runners
or between the extra and retained runners are allowed. Consequently
a 58-runner configuration of radius at most \(1/59\) can contain at
most 50 of the original 57 speed--phase pairs.
\end{theorem}
\begin{proof}
Use \(Q=59000000\) and the coordinate \(x=Qt\). The open gaps
\((a,b)\) in the closed bands of \(B\) at \(d_0\) number 207.
In each take the integer interval
\([\lceil Qa\rceil,\lfloor Qb\rfloor]\) and use its interior.
The companion file \texttt{residual\_weights.json} gives the
nonnegative integer density \(w_j\) on each, in increasing order of
\(a\). The 121 positive-length intervals of positive weight have
total weight \(W=\sum_j w_j=303345\) and total weighted length
\(T=\sum_j w_j(b_j-a_j)=153120431\), where \(a_j,b_j\) are the
integer endpoints. Normalize to probability mass by dividing the
densities by \(T\). For each supported interval, each old runner's
phase on it is between \(k+d_0\) and \(k+1-d_0\) for one integer
\(k\), as a rational endpoint check confirms. Its interior lies
outside all the old bands, and the seed covers measure zero for
this probability measure.

At speed \(v\), change coordinates again to \(y=vx\).
Modulo \(Q\), the covered \(y\)-coordinates lie in a sliding
window of length \(U=2d_0Q=2000000\). The contribution from
\([a_j,b_j]\) consists of
\(\lfloor v(b_j-a_j)/Q\rfloor\) whole periods and one remainder arc.
Each period contributes \(U\) to the intersection length.
The remainder contributes the arc's intersection length with the
sliding window. The sum of these lengths, with weights \(w_j\),
divided by \(vT\), is the covered mass.
For fixed \(v\) the length is continuous and piecewise affine in
the window start. Breakpoints occur where one of the window
endpoints meets a remainder endpoint, which are integral
coordinates modulo \(Q\).

For \(1\le v\le59940\) we enumerate those breakpoints.
A verifier forms the piecewise constant density on the residue
circle by adding the weights of the remainder arcs, integrates it
by a prefix sum of integer rectangle areas, and subtracts the
values at the two window endpoints. Wrapped arcs are split at
zero or handled by periodic extension of the integral.
For each \(v\) we add the independent whole-period term and
compare exact rational masses. The largest over the finite range is
\[
 \beta=\frac{20159589}{153120431},
\]
attained for speed \(160\) and window start \(9192800\).
The second verifier computes the same maximum by accumulating
slope changes of the window intersection function.
For every \(v\ge59941\), the whole periods contribute at most \(U/Q\)
to the normalized mass, and the remainders contribute at most
\(UW/(vT)\). Thus the mass is bounded by
\[
 \frac UQ+\frac{UW}{vT}
 \ \le\ \frac{7735870501306}{77359044788527}
 \ <\ \beta .
\]
Seven new runners cover mass at most
\(7\beta=141117123/153120431<1\).
This proves the obstruction even with weak coverage.
If a 58-runner cover kept at least 51 of the old pairs,
restoring every missing old pair would leave all 57 old pairs and at
most seven additional pairs, a contradiction.
\end{proof}

Linear programming proposed the weights and two exact interval
checks reproduce the bound. With all weights equal the precursor calculation
\texttt{residual\_mass.py} gives a weaker mass cap \(12371/51040\)
and hence at least five extra runners. An earlier direct
two-runner calculation uses just 20 of the seed gaps.

On an interval of length \(L\), a runner of speed \(v\) covers
measure at most \(2d_0L+2d_0/v\), by splitting into full periods and
one remainder. For two runners with speeds \(v\le w\), the longest
gap, of width \(26269/295000000\), forces
\(v\le 4d_0/((1-4d_0)L)\), hence \(v\le816\).
Divide the slower phase into 32 closed boxes of half-width \(1/64\).
A gap with midpoint \(t\) and width \(l\) is missed throughout a box
centred at \(s\) if
\(\norm{vt+s}>d_0+1/64+vl/2\).
For each of the \(26112\) speed-box pairs among the 20 widest gaps
at least eight are missed. The fast speed has finite upper bound
\(2d_0/l\) from the widest such gap, and for every two missed
midpoints \(t_i,t_j\) it must satisfy
\(\norm{w(t_i-t_j)}\le2d_0\).
All candidate sets are empty by rational arithmetic.
A separate single-runner check eliminates the unused speeds through
380 at \(d_0\) using the four widest gap constraints.
The two-runner calculation permits both speeds to coincide.

\section{Deletion and insertion with fixed surviving phases}

We also tested reducing each of two fixed configurations by one
runner. The first seed is the zero-phase speed set \(1,\ldots,59\)
and the second is Poliakova's 95-runner construction, built by
applying rule \(A\) to \(2,4,\ldots,192\) with its phases evaluated
at \(1/97\) \cite[Theorem 7]{Poliakova}. The target radii are
\(1/59\) and \(1/95\), respectively. This section uses strict coverage.
Each new speed must be distinct from the surviving speeds, and the
two new speeds, when present, distinct from each other. A deleted
speed may be re-used with a different phase.

\subsection*{Deleting two and inserting one}

Sweep the endpoints of all original open bands. Record each endpoint
and open cell with its exact set of covering runners whenever that
set has size at most two. At an endpoint first remove arcs ending
there, record the endpoint, and then insert arcs starting there.
Deleting a specified pair exposes exactly the pieces whose covering
sets lie in the pair. The set so exposed is closed, so we can merge
the closures of its cells and include every exposed singleton.
If the widest residual component has positive length \(L\), the new
speed is at most \(\lfloor 2d/L\rfloor\) and its allowable phase
windows are the open form of (1), including the constraints of
singletons. If the residual has only rational singletons, a multiple
of their denominators hits all of them at phase zero. An unused such
speed exists. The empty residual requires no phase constraint.
The resulting complete finite computations give the following.
\begin{center}
\begin{tabular}{@{}lrrr@{}}
\toprule
Target & Deletion pairs & Largest speed bound & Feasible \\
\midrule
58 & 1711 & 170 & 0 \\
94 & 4465 & 580 & 0 \\
\bottomrule
\end{tabular}
\end{center}
See \texttt{single\_replacement.py} and its exact results.

If the exposed compact set is symmetric about time zero and
\(d<1/4\), a strict single-runner cover can be chosen at phase \(0\)
or \(1/2\). Indeed for each exposed pair \(t,-t\), both \(s+vt\) and
\(s-vt\) are within \(d\) of integers \(k,l\), respectively.
The average of the signed errors is \(s-(k+l)/2\) and half their
difference is \(vt-(k-l)/2\). Modulo one there is a unique phase in
\(\{0,1/2\}\) within \(d\) of \(s\), independent of \(t\).
It is the residue of both half-integers \((k+l)/2,(k-l)/2\).
Thus the corresponding phase gives distance less than \(d\) at every
exposed time. The computation above treats the continuous phases
directly in both seeds.

\subsection*{Deleting three and inserting two}

We selected twelve triples per seed by the smallest number of exposed
positive-length cells before merging, breaking ties by seed order.
Sweep all open bands as above, now keeping covering sets of size up
to three. For each triple let \(H\) be the union of closed residual
components. We reduce the search on \(H\) as follows.

Let \(v<w\) be the new speeds and let \(S=2d/v,\ F=2d/w\) be their
band lengths. For \(d<1/4\) any complete gap between two slow
bands is longer than a fast band. Thus each covered component of
\(H\) can meet at most one slow band. Its two portions outside that
band each have length strictly less than \(F\) because each is
contained in one open fast band. If it misses the slow bands
entirely its length is less than \(F\). Therefore the maximum length
\(L\) of a residual component requires
\[
  L<S+2F<6d/v .
\]
This caps \(v\), and caps \(w\) whenever \(L>S\).
When a component \([a,b]\) has length at least \(F\), it must meet
a slow band, whose endpoints must be strictly below \(a+F\)
and strictly above \(b-F\), respectively. This requires the slow
phase to lie in
\[
   (-d-va-vF,\ d-vb+vF)\pmod1 .
\]
Intersect these necessary conditions from all qualifying components.
Failure persists as \(w\) increases, since \(F\) decreases.
Two more finite filters can be used. If \(g=\gcd(v,w)\) and
\(gL\ge1\), then one common period \(1/g\) inside \(H\) cannot be
covered by the two runners, whose combined covered fraction is
at most \(4d<1\). If \(L\ge S+F\), both wings of a longest component
require distinct fast bands. If their centers differ by \(m/w\),
the component's endpoints force \(L<m/w+F\), while both fast bands
meeting the respective slow-band boundaries forces \(m/w-F<S\).

For each remaining fixed speed pair, split the slow and fast phase
axes at the finitely many values where a missed time band starts
or stops meeting a component of \(H\).
For a speed \(v\) and phase \(s\) its closed missed band with
index \(k\) is
\([(k+d-s)/v,(k+1-d-s)/v]\).
The set of relevant indices is constant on the open phase cells.
For a fast phase \(b_0\), a missed time in one residual component
exists precisely when that component and one closed gap from each
runner meet pairwise. Three closed real intervals meet all together
if they meet pairwise. The two gaps with indices \(k,j\) meet
if and only if
\[
 w(k+d)-v(j+1-d)\ \le\ ws-vb_0\ \le\
 w(k+1-d)-v(j+d).
\]
Thus on each open phase rectangle feasibility is the existence of
a point in its open interval of values \(ws-vb_0\) outside a finite
union of closed intervals. For fixed \(k\), consecutive \(j\)'s
translate these forbidden blocks by \(v\). Their common length
\((w+v)(1-2d)\) exceeds \(v\), so they can be merged before the
interval-union test. A cover strict on the compact set \(H\)
has a positive margin and persists under small phase changes.
It is sufficient to test only open phase rectangles, with every
time endpoint still included in \(H\).

For the 24 selected triples the exact solver proved infeasibility.
It used a fallback speed cap of 2500, but for each allowed slower
speed either a finite analytic bound covered the entire fast range
or the monotone phase condition excluded the entire remaining
tail. The solver distinguishes these global exclusions from
cap-limited negative outcomes. The exact solver is
\texttt{two\_speed\_search.py}. The driver \texttt{triple\_replacement.py}
selects the triples and \texttt{triple\_replacement.json} records
each verdict. The following are the triples, listed as speeds.
\begin{center}
\begin{tabular}{@{}cc@{}}
\toprule
Target 58 & Target 94 \\
\midrule
\((1,2,6)\)  & \((2,3,9)\) \\
\((1,2,59)\) & \((3,4,9)\) \\
\((1,3,4)\)  & \((2,4,9)\) \\
\((1,3,59)\) & \((4,6,9)\) \\
\((1,4,6)\)  & \((2,3,4)\) \\
\((1,4,59)\) & \((2,3,6)\) \\
\((1,6,59)\) & \((2,6,9)\) \\
\((2,3,59)\) & \((3,4,6)\) \\
\((2,6,59)\) & \((3,6,9)\) \\
\((1,2,3)\)  & \((2,3,16)\) \\
\((1,2,4)\)  & \((2,9,16)\) \\
\((2,3,4)\)  & \((3,9,16)\) \\
\bottomrule
\end{tabular}
\end{center}

\section{Tests in phase boxes}

For fixed distinct positive speeds and \(0<d<1/2\), the family of
open runner bands covers \(\T\) if and only if every right endpoint
of every band lies in another band. For necessity the band that
ends there is open and no other band of the same runner contains
that endpoint. For sufficiency, a proper finite nonempty union
of open arcs has a boundary point at a right endpoint of one arc,
including the case of a singleton complement next to it.
By compactness pointwise strict coverage is equivalent to
\(\gamma<d\).

Let the nominal phases be \(s_i\) and vary them to \(s_i-v_i y_i\)
with \(|v_i y_i|\le\epsilon\). For each original right endpoint
\(e=(k+d-s_i)/v_i\), the exact condition is that at least one choice
of \(j\ne i\) and \(l\in\Z\) satisfies
\[
 \frac{l-s_j-v_j e-d}{v_j}< y_i-y_j <
 \frac{l-s_j-v_j e+d}{v_j}.
\]
The bounds on the delays exclude all but finitely many \(l\).
After picking a disjunct for each endpoint, this is a system
of strict and weak difference inequalities with a zero reference
for the delay box.

For rational nominal phases, radius and half-width let \(D\) be
the product of their common denominator and the lcm of the speeds.
All bounds after scaling delays by \(D\) are integral.
A real feasible selection has nonnegative weight on every
directed cycle and strictly positive weight on each cycle using
a strict edge. Every simple cycle has at most \(n+1\) edges, so
tightening each strict bound by \(1/(n+1)\) leaves all cycle weights
nonnegative. Multiply again by \(n+1\).
Integer shortest-path potentials then give a point satisfying the
tightened constraints. Subtract the potential of the reference
from every coordinate to reset it to zero. Conversely any point on
this grid satisfying those tightened constraints solves the original
strict ones. This reduction is exact for continuous phases and time.
The interval and difference-constraint approach is related to
\cite[Section 5]{BCS}.

At \(d=1/59\) the exact integer-difference encoding in Z3 returned
infeasible in the following two boxes.
\begin{center}
\begin{tabular}{@{}lrr@{}}
\toprule
Fixed speeds and nominal phases & Half-width & Endpoints \\
\midrule
\(1,\ldots,58\), zero phases & \(3/125\) & 1711 \\
Stored 58-runner candidate & \(1/32\) & 3217 \\
\bottomrule
\end{tabular}
\end{center}
A separate real-arithmetic Z3 encoding independently excludes the
standard-speed case at half-width \(1/50\) using 203 endpoints.
The two larger results use solver verdicts from the integer encoding.
The real encoding timed out on those larger boxes.
The nominal phases for the second row are the exact rationals in
\texttt{chain58\_box.json}, obtained from the stored search output.
A common time translation of all phases preserves the conclusion for
each box.

An additional finite-menu SAT calculation used the full time circle
at \(64/(64K+1)\), for \(K=59\), with at most \(K-1\) runners
chosen from finite speed and phase options. The union of closed
coverage arcs covers the circle precisely when it covers every open
cell cut out by their rational endpoints. A clause requires at least
one active option per cell not covered by a fixed runner. Further
constraints allow at most one phase per speed and enforce the
runner budget. Seven menus at \(K=59\) were infeasible at that fixed
radius. The exact parameters and verdicts are in
\texttt{direct\_rule\_search.log} and
\texttt{direct\_rule\_frontier.log}. The radius is strictly smaller
than \(1/K\), so even within a searched menu the verdict is scoped to
that fixed radius. A global 58-runner question remains.

\section*{Data, formalization and attribution}

The exact data and replay commands are in the
\href{https://github.com/F1Justin/MathLab/tree/v1.0.0}{versioned MathLab repository}.
The \texttt{MathLab.Frozen} namespace contains
\texttt{frozen\_optimum}, its attaining cover, and
\texttt{perturbation\_barrier}. The formalization
concerns extending and perturbing the already published seed.
The seed and its exact radius are credited to \texttt{tomdif};
\texttt{aadrianvv} provided the earlier certificates and
the upstream public verifier. The upstream Lean cover checker,
which was reused under its MIT license, has its separate attribution
in the formal project. The finite phase and replacement checks
use rational or integer arithmetic as described above.

Chengyi Fang initiated and funded the project, set its objectives,
supervised its execution and directed dissemination. OpenAI Codex led
proof development, computation, formalization and integration. GPT-6.2 Sol
subagents assisted with mathematical auditing and publication. GPT-6.2 Astra
developed the exact two-runner phase solver and assisted with formal proofs
and independent model review through the Responses API. AI systems produced
the mathematics, software and manuscript. Chengyi Fang has not independently
checked any part of the proofs.

\medskip\noindent
This work is licensed under
\href{https://creativecommons.org/licenses/by/4.0/}{CC BY 4.0}.
The original accompanying code is released under the MIT License.

\begingroup
\small
\begin{thebibliography}{9}
\setlength{\itemsep}{0pt}
\bibitem{BCS}
M.~Blanco, F.~Criado and F.~Santos,
\emph{Coloopless zonotopes and counterexamples to the Shifted Lonely
Runner Conjecture}, \href{https://arxiv.org/abs/2603.24784v4}
{arXiv:2603.24784v4}, 2026.
\bibitem{Poliakova}
D.~Poliakova, \emph{More (shifted) runners, less loneliness},
\href{https://arxiv.org/abs/2609.23952v1}{arXiv:2609.23952v1}, 2026.
\bibitem{Zhang}
Y.~Zhang, \emph{Tight instances of the Lonely Runner Conjecture:
one-entry modifications and the growth of speeds},
\href{https://arxiv.org/abs/2608.13599v3}{arXiv:2608.13599v3}, 2026.
\bibitem{upstream}
\texttt{aadrianvv} and \texttt{tomdif},
\emph{Shifted Lonely Runner counterexamples for \(18\le n\le57\)},
\href{https://github.com/1337isnot1337/solved-conjectures/issues/110}
{solved-conjectures, issue 110}, 2026.
\bibitem{upstreamlean}
\texttt{1337isnot1337/solved-conjectures},
\emph{Lean-verify the shifted Lonely Runner counterexamples for 18 through 57},
\href{https://github.com/1337isnot1337/solved-conjectures/commit/9aefed33ac284780cf29b76d7914bc32877757af}
{pull request 755, merged 26 November 2026}.
\end{thebibliography}
\endgroup
\end{document}
